\section{命令实现示例}

\subsection{\texttt{cmd\_cls}：清屏}

\codefile{src/kernel/cmds/cmd\_cls.c}
\begin{lstlisting}[style=cstyle]
printk("\x1b[2J\x1b[H");
\end{lstlisting}

\begin{itemize}
  \item \texttt{\textbackslash x1b[2J}：ANSI 转义序列——清除整个屏幕
  \item \texttt{\textbackslash x1b[H}：光标移到左上角 (1,1)
\end{itemize}

这只对支持 ANSI 的串口终端有效。QEMU \texttt{-serial stdio} 的宿主终端通常支持。

\subsection{\texttt{cmd\_shutdown}：关机}

\codefile{src/kernel/cmds/cmd\_shutdown.c}
\begin{lstlisting}[style=cstyle]
outw(0x604, 0x2000);   // QEMU ACPI 关机
outw(0xB004, 0x2000);  // Bochs 关机
cpu_halt_forever();     // cli + hlt 循环
\end{lstlisting}

\begin{whybox}
为什么要尝试两个不同的端口？因为 QEMU 和 Bochs 的 ACPI 关机端口不同。
如果两个都无效（真实硬件），就退化为 \texttt{cli; hlt}——CPU 停止，
电源不会关闭但系统不再响应。
\end{whybox}

% ============================================================================

